\section{Qué «oído» aprendió el modelo: Inspecting the Learned Front End}

El mAP por sí solo todavía no dice cómo representa el modelo el sonido. El artículo examina además los filters de la primera capa: a cada learned filter se le aplica la Fourier transform, se ordenan por center frequency y luego se observan la densidad de bandas de frecuencia y la filter shape. Este proceso convierte los parámetros de la red neuronal en respuestas en frecuencia, que los lectores de procesamiento de señales ya conocen, y es la parte de esta clase más cercana a la mechanistic inspection.

\subsection{De la Filter Impulse Response a la Frequency Response}

Si los pesos en el dominio del tiempo del $j$-ésimo filter del front end son $h_j[n]$, su respuesta en frecuencia es
\[
H_j(\omega)=\sum_n h_j[n]e^{-i\omega n}.
\]
La posición del pico permite definir una center frequency aproximada
\[
\omega_j^*=\operatorname*{argmax}_{\omega}|H_j(\omega)|.
\]
Al ordenar los filters según $\omega_j^*$, un Fourier bank fijo se distribuye de manera uniforme, aproximadamente a lo largo de una recta; si el learned bank asigna más filters a ciertas bandas de frecuencia, la curva se dobla o presenta un step pattern, lo que indica que el modelo redistribuye la resolution en frecuencia según la tarea.

\lecturefigure{slide-45-ideal-time-frequency.jpg}{Comparación entre un Fourier bank fijo ideal y una learned task-specific frequency organization.}{Video oficial de Stanford Online, 00:44:36--00:46:09.}

\begin{knowledgebox}{Lectura de la figura: qué representan el eje horizontal, el eje vertical y la curva}
El eje horizontal es el filter index después del ordenamiento; el eje vertical es la center frequency en la que cada filter tiene su respuesta máxima. Una recta indica una distribución uniforme de frecuencias; una curva indica que ciertas bandas reciben filters más densos. Que tareas distintas produzcan curvas distintas respalda que «el front end se adapta a la tarea», pero la curva por sí misma no dice qué representación es más robusta fuera de la distribución.
\end{knowledgebox}

\teachervoice{Verma describe este fenómeno como una computadora que «ajusta sus propios oídos» según la tarea. Esta frase es más precisa que «el modelo aprendió un spectrogram»: lo que aprende es un conjunto de analysis filters no uniformes y dependientes de la tarea, no una copia de una STFT fija.}

\subsection{DSP-like Filters: Onset, Window y Sinusoid}

Al visualizar los individual filters se observan formas semejantes a un onset detector, a una windowing function y a una sinusoidal basis; en una misma banda de frecuencia también pueden aparecer varios filters con fases distintas. Esto muestra que el gradient descent redescubrió varios building blocks clásicos de DSP y, al mismo tiempo, combinó formas que una basis fija no prevé.

\lecturefigure{slide-46-learned-filter-bank.jpg}{Learned filters: patrones onset-sensitive, window-like, sinusoidal y phase-diverse.}{Video oficial de Stanford Online, 00:46:10--00:47:18.}

\begin{warningbox}{La visual resemblance no es una prueba completa del mecanismo}
Que un filter se parezca a una Hamming window o a una sinusoid solo indica que la shape/response es similar. Para demostrar una causal function se requieren además intervention, ablation, pruebas con entradas sintetizadas y estabilidad entre checkpoints; varios filters también pueden implementar en conjunto una distributed feature.
\end{warningbox}

\begin{importantbox}{Relación entre el DSP tradicional y el End-to-End Learning}
La conclusión más valiosa de esta clase no es que «las redes neuronales sustituyen al procesamiento de señales», sino que ambos pueden formar un ciclo cerrado: el DSP aporta bases interpretables, multi-scale transforms y herramientas de diagnóstico; el end-to-end learning redistribuye las bandas de frecuencia y combina filters según la tarea; y los resultados del análisis orientan, a su vez, el diseño de la architecture.
\end{importantbox}

\subsection{Resumen de la sección}

La Fourier inspection proyecta los learned weights sobre la center frequency y la filter shape, y muestra que el front end forma un filterbank dependiente de la tarea y reproduce parte de la DSP structure clásica. Es una evidencia sólida sobre la representación, pero todavía falta un paso para llegar a un causal circuit completo.
